\chapter{Essential Physical Chemistry: Complete Numerical Bank}
\label{chap:essential_all_questions}

\begin{problembox}
\textbf{Problem 1} \hfill \textbf{[Essential Physical Chemistry p.42 Q1]}
\label{q:ext-essential_all_questions-1}

Define or explain the following terms :
(a)
Arrhenius theory
(b)
Transport number
(c)
Hittorf's method
(d)
Moving boundary method
(e)
Kohlrausch's law
(f)
Degree of dissociation
\end{problembox}

\begin{problembox}
\textbf{Problem 2} \hfill \textbf{[Essential Physical Chemistry p.42 Q2]}
\label{q:ext-essential_all_questions-2}

(a)
Explain ionic conductance, transport number of an ion, and absolute ionic velocity.
(b)
The absolute velocity of Ag+ is 0.00057 cm sec--1 and of the NO3
-- is 0.00069 cm sec--1. Assuming
complete dissociation, calculate the specific conductivity of 0.01 M solution of silver nitrate.
Answer. (b) 0.00012 mhos
\end{problembox}

\begin{problembox}
\textbf{Problem 3} \hfill \textbf{[Essential Physical Chemistry p.42 Q3]}
\label{q:ext-essential_all_questions-3}

Describe Hittorf's method for the experimental determination of the transference number. The speed
ratio of silver and nitrate ions in AgNO3 electrolysed between silver electrodes was found to be 0.89.
Calculate the transference number of silver and nitrate ions.
Answer. Ag+ = 0.4708; NO3
-- = 0.5292
\end{problembox}

\begin{problembox}
\textbf{Problem 4} \hfill \textbf{[Essential Physical Chemistry p.42 Q4]}
\label{q:ext-essential_all_questions-4}

(a)
What do you understand by the transport number of an ion? Does it vary?
(b)
In a moving boundary experiment with 0.1 N KCl using 0.65 N LiCl as indicator solution, a
constant current of 0.006 amp was passed for 1900 secs and the boundary was observed to move
through 5 cm in a tube of 0.1142 cm2 cross-section. Calculate the transport number of K+ and Cl--
ions.
Answer. (b) K+ = 0.4833; Cl-- = 0.5167
\end{problembox}

\begin{problembox}
\textbf{Problem 5} \hfill \textbf{[Essential Physical Chemistry p.42 Q5]}
\label{q:ext-essential_all_questions-5}

(a)
What is meant by transport number of ions? Describe Hittorf's method for the determination of
transport number of silver ions.
(b)
A dilute solution of silver nitrate was electrolysed using platinum electrodes. After passing small
\end{problembox}

\begin{problembox}
\textbf{Problem 6} \hfill \textbf{[Essential Physical Chemistry p.42 Q1]}
\label{q:ext-essential_all_questions-6}

Conductance measurements are done to
check end points.
\end{problembox}

\begin{problembox}
\textbf{Problem 7} \hfill \textbf{[Essential Physical Chemistry p.42 Q2]}
\label{q:ext-essential_all_questions-7}

Titrations can be carried out even with
coloured solution.
\end{problembox}

\begin{problembox}
\textbf{Problem 8} \hfill \textbf{[Essential Physical Chemistry p.42 Q3]}
\label{q:ext-essential_all_questions-8}

Accurate results are obtained.
\end{problembox}

\begin{problembox}
\textbf{Problem 9} \hfill \textbf{[Essential Physical Chemistry p.42 Q4]}
\label{q:ext-essential_all_questions-9}

End point is determined graphically.
\end{problembox}

\begin{problembox}
\textbf{Problem 10} \hfill \textbf{[Essential Physical Chemistry p.42 Q5]}
\label{q:ext-essential_all_questions-10}

In case of polybasic acids conductometric
titrations can be used.
\end{problembox}

\begin{problembox}
\textbf{Problem 11} \hfill \textbf{[Essential Physical Chemistry p.42 Q6]}
\label{q:ext-essential_all_questions-11}

These are successful even in weak acids
and strong bases.
\end{problembox}

\begin{problembox}
\textbf{Problem 12} \hfill \textbf{[Essential Physical Chemistry p.42 Q1]}
\label{q:ext-essential_all_questions-12}

Volume measurements are done to check
end points.
\end{problembox}

\begin{problembox}
\textbf{Problem 13} \hfill \textbf{[Essential Physical Chemistry p.42 Q2]}
\label{q:ext-essential_all_questions-13}

These titrations fail  in coloured solutions
as suitable indicators are not available some
times.
\end{problembox}

\begin{problembox}
\textbf{Problem 14} \hfill \textbf{[Essential Physical Chemistry p.42 Q3]}
\label{q:ext-essential_all_questions-14}

Results are not so accurate.
\end{problembox}

\begin{problembox}
\textbf{Problem 15} \hfill \textbf{[Essential Physical Chemistry p.42 Q4]}
\label{q:ext-essential_all_questions-15}

End point is determined by change in
colour of indicator.
\end{problembox}

\begin{problembox}
\textbf{Problem 16} \hfill \textbf{[Essential Physical Chemistry p.42 Q5]}
\label{q:ext-essential_all_questions-16}

In case of polybasic acids volumetric
titrations do not give correct end points.
\end{problembox}

\begin{problembox}
\textbf{Problem 17} \hfill \textbf{[Essential Physical Chemistry p.42 Q6]}
\label{q:ext-essential_all_questions-17}

These are not successful in weak acids and
strong bases.
\end{problembox}

\begin{problembox}
\textbf{Problem 18} \hfill \textbf{[Essential Physical Chemistry p.43 Q6]}
\label{q:ext-essential_all_questions-18}

(a)
Explain Kohlrausch's law of ionic mobility. How does it help in determining the equivalent
conductivity of infinite dilution of weak electrolyte?
(b)
Describe the Hittorf's method for determining the transport number of Ag+ and NO3
-- ions in
solution of silver nitrate.
\end{problembox}

\begin{problembox}
\textbf{Problem 19} \hfill \textbf{[Essential Physical Chemistry p.43 Q7]}
\label{q:ext-essential_all_questions-19}

The specific conductance of a saturated solution of AgCl is 1.55 $/times$ 10--6 $/Omega$--1 cm--1 at 25$/circ$C. The ion
conductance of Ag+ Cl--1 are 61.94 and 76.34 $/Omega$--1 cm2 respectively. Calculate the solubility of AgCl in
gram equivalent per litre at 25$/circ$C. Neglect the specific conductance of water.
Answer. 1.6148 $/times$ 10--3 g lit--1
\end{problembox}

\begin{problembox}
\textbf{Problem 20} \hfill \textbf{[Essential Physical Chemistry p.43 Q8]}
\label{q:ext-essential_all_questions-20}

In a moving boundary experiment with 0.1 N KCl solution, the boundary moved 4.94 cm during
67 minutes when a current of 5.21 A was used. The cross-sectional area of the tube was 0.230 cm2.
Calculate the transport number of K+ ions.
Answer. 0.523
\end{problembox}

\begin{problembox}
\textbf{Problem 21} \hfill \textbf{[Essential Physical Chemistry p.43 Q9]}
\label{q:ext-essential_all_questions-21}

What do you understand by the terms: transport number and ionic mobility of an ion? How are these
related to each other? Write their units.
(Punjabi BSc, 2000)
\end{problembox}

\begin{problembox}
\textbf{Problem 22} \hfill \textbf{[Essential Physical Chemistry p.43 Q10]}
\label{q:ext-essential_all_questions-22}

What do you understand by transport number? Give one method for the determination of it.
(Punjabi BSc, 2000)
\end{problembox}

\begin{problembox}
\textbf{Problem 23} \hfill \textbf{[Essential Physical Chemistry p.43 Q11]}
\label{q:ext-essential_all_questions-23}

Describe Hittorf's method for determining the transport number of Ag+ and NO3
-- ions in a solution of
AgNO3 when the electrodes are made up of silver.
(Jiwaji BSc, 2000)
\end{problembox}

\begin{problembox}
\textbf{Problem 24} \hfill \textbf{[Essential Physical Chemistry p.43 Q12]}
\label{q:ext-essential_all_questions-24}

Give the applications of transport number.
(Himachal Pradesh BSc, 2000)
\end{problembox}

\begin{problembox}
\textbf{Problem 25} \hfill \textbf{[Essential Physical Chemistry p.43 Q13]}
\label{q:ext-essential_all_questions-25}

State Kohlrausch's law.
(Madurai BSc, 2000)
\end{problembox}

\begin{problembox}
\textbf{Problem 26} \hfill \textbf{[Essential Physical Chemistry p.43 Q14]}
\label{q:ext-essential_all_questions-26}

Write short notes on :
(a)
Transport number
(b)
Kohlrausch law of independent migration of ions (Delhi BSc, 2000)
\end{problembox}

\begin{problembox}
\textbf{Problem 27} \hfill \textbf{[Essential Physical Chemistry p.43 Q15]}
\label{q:ext-essential_all_questions-27}

(a)
How is solubility of sparingly soluble salt is determined by conductance measurement?
(b)
Discuss the application of Kohlrausch's law.
(Lucknow BSc, 2001)
\end{problembox}

\begin{problembox}
\textbf{Problem 28} \hfill \textbf{[Essential Physical Chemistry p.43 Q16]}
\label{q:ext-essential_all_questions-28}

If the conductance at infinite dilution of NaCl, HCl and CH3COONa are 125.45 $/times$ 10--4, 426.1 $/times$ 10--4 and
91.0 $/times$ 10--4 cm2 eqvt--1 respectively, calculate equivalent conductance of acetic acid at infinite dilution.
Answer. 390.7 $/times$ 10--4 $/Omega$--1 cm2 eqvt--1
(Guru Nanak Dev BSc, 2002)
\end{problembox}

\begin{problembox}
\textbf{Problem 29} \hfill \textbf{[Essential Physical Chemistry p.43 Q17]}
\label{q:ext-essential_all_questions-29}

Define Transport number. Describe briefly the principle of experimental determination of Transport
number by Hittorf's method.
(Panjab BSc, 2002)
\end{problembox}

\begin{problembox}
\textbf{Problem 30} \hfill \textbf{[Essential Physical Chemistry p.43 Q18]}
\label{q:ext-essential_all_questions-30}

(a)
State the principle of conductometric titrations. Draw the titration curves for
(i)
Weak acid with a strong base
(ii)
Na2SO4 solution with BaCl2 solution
(b)
Calculate the ionisation constant of water if the specific conductivity of water is
5.51$/times$10--8 $/Omega$--1 cm--1 and ionic conductance of H+ and OH-- ions at the same temperature are
349.8 and 196.5 respectively.
Answer. 1.013 $/times$ 10--14
(Arunachal BSc, 2002)
\end{problembox}

\begin{problembox}
\textbf{Problem 31} \hfill \textbf{[Essential Physical Chemistry p.43 Q19]}
\label{q:ext-essential_all_questions-31}

How does Kohlrausch's law helps in determining the dissociation constants of a weak acid by
conductance measurements?
(Vidyasagar BSc, 2002)
\end{problembox}

\begin{problembox}
\textbf{Problem 32} \hfill \textbf{[Essential Physical Chemistry p.43 Q20]}
\label{q:ext-essential_all_questions-32}

Derive the expression for the determination of Transport number by Hittorf's method when electrodes
are not attacked?
(Arunachal BSc, 2002)
\end{problembox}

\begin{problembox}
\textbf{Problem 33} \hfill \textbf{[Essential Physical Chemistry p.43 Q21]}
\label{q:ext-essential_all_questions-33}

Explain the factors on which the transport number depends.
(MD Rohtak BSc, 2002)
\end{problembox}

\begin{problembox}
\textbf{Problem 34} \hfill \textbf{[Essential Physical Chemistry p.43 Q22]}
\label{q:ext-essential_all_questions-34}

Explain why the transport number of cadmium ion in cadmium iodide solution at high concentration
may become zero or negative also.
(Arunachal BSc, 2002)
\end{problembox}

\begin{problembox}
\textbf{Problem 35} \hfill \textbf{[Essential Physical Chemistry p.43 Q23]}
\label{q:ext-essential_all_questions-35}

Explain how is Kohlrausch's law helpful in determining the equivalent conductance of weak electrolyte
at infinite dilution.
(Jammu BSc, 2002)
\end{problembox}

\begin{problembox}
\textbf{Problem 36} \hfill \textbf{[Essential Physical Chemistry p.44 Q24]}
\label{q:ext-essential_all_questions-36}

Write short notes on :
(i)
Activity and mean ionic activity of an electrolyte.
(ii)
Ionic atmosphere
(Aligarh BSc, 2002)
\end{problembox}

\begin{problembox}
\textbf{Problem 37} \hfill \textbf{[Essential Physical Chemistry p.44 Q25]}
\label{q:ext-essential_all_questions-37}

Discuss the effect of concentration on transference number.
(Kalyani BSc, 2003)
\end{problembox}

\begin{problembox}
\textbf{Problem 38} \hfill \textbf{[Essential Physical Chemistry p.44 Q26]}
\label{q:ext-essential_all_questions-38}

During the electrolysis of a solution of potassium chloride between platinum electrodes, 0.0137 g of
chloride was lost from the anodic compartment and 0.0857 g of silver was deposited in a silver
coulometer connected in series with the cell. Find at the transport number of K+ and Cl-- ions.
Answer. 0.4854; 0.5146
(Delhi BSc, 2003)
\end{problembox}

\begin{problembox}
\textbf{Problem 39} \hfill \textbf{[Essential Physical Chemistry p.44 Q27]}
\label{q:ext-essential_all_questions-39}

(a)
Explain the term 'ionic mobility'. Describe Hittorf's theoretical device to show that although most
of the ions differ largely in their mobilities, but their equivalent amounts are discharged on electrolysis
at appropriate electrodes.
(b)
In moving boundary experiment with 0.100 N KCl using 0.065 N LiCl as indicator solution, a
constant current of 0.005893 amp was passed for 2130 seconds. The boundary was observed to
move through 5.6 cm in a tube of 0.1142 sq cm cross-section. Calculate the transport number of
K+ and Cl-- ions.
Answer. 0.492; 0.508
(Rajasthan BSc, 2003)
\end{problembox}

\begin{problembox}
\textbf{Problem 40} \hfill \textbf{[Essential Physical Chemistry p.44 Q28]}
\label{q:ext-essential_all_questions-40}

What is meant by limiting equivalent conductance of weak electrolyte? How do you determine it by the
help of Kohlrausch's law?
(Sambalpur BSc, 2003)
\end{problembox}

\begin{problembox}
\textbf{Problem 41} \hfill \textbf{[Essential Physical Chemistry p.44 Q29]}
\label{q:ext-essential_all_questions-41}

Explain the use of Kohlrausch's law to determine the solubility of a sparingly soluble salt.
(Nagpur BSc, 2003)
\end{problembox}

\begin{problembox}
\textbf{Problem 42} \hfill \textbf{[Essential Physical Chemistry p.44 Q30]}
\label{q:ext-essential_all_questions-42}

(a)
Define transport number.
(b)
A solution of silver nitrate containing 0.51 gm of the salt in 60.40 gm of the solution was electrolysed
under similar electrodes. After electrolysis 64 gm of the anode solution was found to contain 0.56
gm of silver nitrate. A current of 0.04 amp was passed for 2950 seconds. Calculate the transference
number of nitrate ion.
Answer. (b) 0.520
(Delhi BSc, 2004)
\end{problembox}

\begin{problembox}
\textbf{Problem 43} \hfill \textbf{[Essential Physical Chemistry p.44 Q31]}
\label{q:ext-essential_all_questions-43}

(a)
How is Kohlrausch's law helpful in determining the solubility of a sparingly soluble salt in water?
(b)
How can you find the transference number of an ion by Hittorf's method?
(Madurai BSc, 2004)
\end{problembox}

\begin{problembox}
\textbf{Problem 44} \hfill \textbf{[Essential Physical Chemistry p.44 Q32]}
\label{q:ext-essential_all_questions-44}

A solution of HCl was electrolysed in a cell between two platinum electrodes. The cathode compartment
contained 0.354 g of chloride ions before electrolysis and 0.32 g after electrolysis. A silver coulometer in
series had a deposit of 0.532 g of silver after passing the same amount of current. Calculate the transport
number of H+ and Cl-- ions.
Answer. 0.84 ;  0.160
(Rajasthan BSc, 2005)
\end{problembox}

\begin{problembox}
\textbf{Problem 45} \hfill \textbf{[Essential Physical Chemistry p.44 Q33]}
\label{q:ext-essential_all_questions-45}

In the Hittorf cell using silver electrodes and AgNO3 as an electrolyte, a certain amount of current was
passing which  deposited  9.886  $/times$  10--4 g equivalent of silver in the coulometer. The anode compartment
had the composition 28.235 g of H2O and 0.099 g of AgNO3 before electrolysis and 28.435 g of water and
0.1874 g of AgNO3 after electrolysis. Calculate the transport number of Ag+ and NO3
-- ions.
Answer. 0.48 and 0.52
(Kalyani BSc, 2005)
\end{problembox}

\begin{problembox}
\textbf{Problem 46} \hfill \textbf{[Essential Physical Chemistry p.44 Q34]}
\label{q:ext-essential_all_questions-46}

In a moving boundary experiment with 1 N KCl solution using CaCl2 as indicator, a current of 0.0115 A
was passed for half an hour and the boundary moved through a volume of 0.106 ml.  Calculate the
transport number of K+ ion.
Answer. 0.494
(Vikram BSc, 2006)
\end{problembox}

\begin{problembox}
\textbf{Problem 47} \hfill \textbf{[Essential Physical Chemistry p.44 Q35]}
\label{q:ext-essential_all_questions-47}

A conductivity cell whose cell constant is 2 cm--1, is filled with 0.1 M acetic acid solution. Its resistance
is found to be 3765 $/Omega$--1. Calculate the degree of dissociation of acetic acid.
Answer. 0.0136
(Delhi BSc, 2006)
\end{problembox}

\begin{problembox}
\textbf{Problem 48} \hfill \textbf{[Essential Physical Chemistry p.44 Q36]}
\label{q:ext-essential_all_questions-48}

In a moving boundary method, a current of 30 m A was passed for 80 seconds. During this time the
boundary of a solution of HCl containing 0.01 mol dm--3 moved 17.0 cm towards the cathode. Calculate
the transport number of H+ and Cl-- ions. (The cross-sectional area of the glass tube is 1.0  $/times$  10--5 m2).
Answer. 0.68 ;  0.32
(Mysore BSc, 2006)
\end{problembox}

\begin{problembox}
\textbf{Problem 49} \hfill \textbf{[Essential Physical Chemistry p.45 Q1]}
\label{q:ext-essential_all_questions-49}

According to Arrhenius theory an electrolyte when dissolved in water gives two types of
(a)
charged particles
(b)
molecules
(c)
ion pairs
(d)
fundamental particles
Answer. (a)
\end{problembox}

\begin{problembox}
\textbf{Problem 50} \hfill \textbf{[Essential Physical Chemistry p.45 Q2]}
\label{q:ext-essential_all_questions-50}

For an electrolyte AB, the ionisation constant is given by the expression
(a)
K = [A] [B]
[AB]
(b)
K = [A ] [B ]
[AB]
+
+
+
(c)
K = [A ] [B ]
[AB]
+
+
-
(d)
K = [A ] [B ]
[AB]
+
-
Answer. (d)
\end{problembox}

\begin{problembox}
\textbf{Problem 51} \hfill \textbf{[Essential Physical Chemistry p.45 Q3]}
\label{q:ext-essential_all_questions-51}

The conductivity of an electrolyte is due to the
(a)
presence of ions in the electrolyte
(b)
free movement of ions in the solution
(c)
reunion of ions in the solution
(d)
release of heat energy due to ionisation
Answer. (b)
\end{problembox}

\begin{problembox}
\textbf{Problem 52} \hfill \textbf{[Essential Physical Chemistry p.45 Q4]}
\label{q:ext-essential_all_questions-52}

The electrical conductivity of an electrolyte depends upon
(a)
the number of molecules in the electrolyte
(b)
the number of ions present in the electrolyte
(c)
the number of ions present in the solution
(d)
the number of molecules of the solvent
Answer. (c)
\end{problembox}

\begin{problembox}
\textbf{Problem 53} \hfill \textbf{[Essential Physical Chemistry p.45 Q5]}
\label{q:ext-essential_all_questions-53}

On passing electrical current through an electrolyte solution
(a)
cations move towards anode
(b)
anions move towards cathode
(c)
cations move towards cathode and anions towards anode
(d)
both cations and anions move in same direction
Answer. (c)
\end{problembox}

\begin{problembox}
\textbf{Problem 54} \hfill \textbf{[Essential Physical Chemistry p.45 Q6]}
\label{q:ext-essential_all_questions-54}

On passing electrical current through an electrolyte solution, the cations
(a)
move towards cathode with the speed equal to that of anions towards anode
(b)
move with faster speed than that of anions
(c)
move with different speed as compared to that of anions
(d)
move with slower speed than that of anions
Answer. (c)
\end{problembox}

\begin{problembox}
\textbf{Problem 55} \hfill \textbf{[Essential Physical Chemistry p.45 Q7]}
\label{q:ext-essential_all_questions-55}

The Hittorf's rule states that
(a)
the loss of concentration around any electrode is proportional to the speed of the ions moving
towards it
(b)
the loss of concentration around any electrode is proportional to the speed of the ions moving
away from it
(c)
the loss of concentrations around both the electrodes is proportional to the sum of speed of
cations and anions
(d)
none of the above
Answer. (b)
\end{problembox}

\begin{problembox}
\textbf{Problem 56} \hfill \textbf{[Essential Physical Chemistry p.45 Q8]}
\label{q:ext-essential_all_questions-56}

The Hittorf's rule can be represented by the expression
\end{problembox}

\begin{problembox}
\textbf{Problem 57} \hfill \textbf{[Essential Physical Chemistry p.46 Q9]}
\label{q:ext-essential_all_questions-57}

The fraction of the total current carried by the cation or anion is termed as
(a)
fractional number
(b)
speed number
(c)
carrier number
(d)
transport number
Answer. (d)
\end{problembox}

\begin{problembox}
\textbf{Problem 58} \hfill \textbf{[Essential Physical Chemistry p.46 Q10]}
\label{q:ext-essential_all_questions-58}

If v+ is the speed of cation and v-- that of anion, the transport number of cation is given by
(a)
+
+
-
$/nu$
$/nu$ -$/nu$
(b)
+
-
+
$/nu$
$/nu$ -$/nu$
(c)
+
+
-
$/nu$
$/nu$ + $/nu$
(d)
-
+
-
$/nu$
$/nu$ + $/nu$
Answer. (c)
\end{problembox}

\begin{problembox}
\textbf{Problem 59} \hfill \textbf{[Essential Physical Chemistry p.46 Q11]}
\label{q:ext-essential_all_questions-59}

If v+ is the speed of cation and v-- that of anion, then the speed ratio is given by
(a)
t
r
t
+
-
=
(b)
1
t
r
t
+
+
=
+
(c)
1
t
r
t
-
+
=
-
(d)
1
t
r
t
+
-
=
-
Answer. (d)
\end{problembox}

\begin{problembox}
\textbf{Problem 60} \hfill \textbf{[Essential Physical Chemistry p.46 Q12]}
\label{q:ext-essential_all_questions-60}

The sum of the transport number of cation and anion is equal to
(a)
1
(b)
0
(c)
0.5
(d)
$/propto$
Answer. (a)
\end{problembox}

\begin{problembox}
\textbf{Problem 61} \hfill \textbf{[Essential Physical Chemistry p.46 Q13]}
\label{q:ext-essential_all_questions-61}

In Hittorf method for determination of transport numbers we make use of a
(a)
H-tube
(b)
V-tube
(c)
U-tube
(d)
L-tube
Answer. (c)
\end{problembox}

\begin{problembox}
\textbf{Problem 62} \hfill \textbf{[Essential Physical Chemistry p.46 Q14]}
\label{q:ext-essential_all_questions-62}

While determining transport number by Hittorf method, \underline{\hspace{1cm}} electrodes are used when electrodes
are attackable
(a)
silver
(b)
copper
(c)
mercury
(d)
platinum
Answer. (d)
\end{problembox}

\begin{problembox}
\textbf{Problem 63} \hfill \textbf{[Essential Physical Chemistry p.46 Q15]}
\label{q:ext-essential_all_questions-63}

In the moving boundary method for the determination of transport number, the formula used is (where
n is the number of faradays of current passed)
(a)
A
sl c
t
n
+ =
(b)
2
A
sl c
t
n
+ =
(c)
2
A
sl c
t
n
+ =
(d)
2
A
sl c
t
n
+ =
Answer. (a)
\end{problembox}

\begin{problembox}
\textbf{Problem 64} \hfill \textbf{[Essential Physical Chemistry p.46 Q16]}
\label{q:ext-essential_all_questions-64}

One Faraday is equal to
(a)
96.500 coulombs
(b)
9650 coulombs
\end{problembox}

\begin{problembox}
\textbf{Problem 65} \hfill \textbf{[Essential Physical Chemistry p.47 Q17]}
\label{q:ext-essential_all_questions-65}

The statement of Kohlrausch's law is
(a)
the equivalent conductance of an electrolyte at infinite dilution is equal to the product of equivalent
conductances of the component ions
(b)
the equivalent conductance of an electrolyte at infinite dilution is equal to the difference of
equivalent conductances of the component ions
(c)
the equivalent conductance of an electrolyte at infinite dilution is equal to the sum of equivalent
conductances of the component ions
(d)
none of the above
Answer. (c)
\end{problembox}

\begin{problembox}
\textbf{Problem 66} \hfill \textbf{[Essential Physical Chemistry p.47 Q18]}
\label{q:ext-essential_all_questions-66}

Kohlrausch's law can be expressed as
(a)
$/lambda$$/alpha$ = $/lambda$a -- $/lambda$c
(b)
$/lambda$$/alpha$ = $/lambda$c -- $/lambda$a
(c)
$/lambda$$/alpha$ = $/lambda$a + $/lambda$c
(d)
$/lambda$$/alpha$ = $/lambda$a $/times$ $/lambda$c
Answer. (c)
\end{problembox}

\begin{problembox}
\textbf{Problem 67} \hfill \textbf{[Essential Physical Chemistry p.47 Q19]}
\label{q:ext-essential_all_questions-67}

Kohlrausch's law can be used to determine
(a)
$/lambda$$/alpha$ for weak electrolytes
(b)
absolute ionic mobilities
(c)
solubility of sparingly soluble salts
(d)
all of these
Answer. (d)
\end{problembox}

\begin{problembox}
\textbf{Problem 68} \hfill \textbf{[Essential Physical Chemistry p.47 Q20]}
\label{q:ext-essential_all_questions-68}

If $/lambda$$/alpha$ is the equivalent conductance at infinite dilution and $/lambda$v the equivalent conductance of the electrolyte
at the dilution v, the degree of dissociation is given by
(a)
v
$/alpha$
$/lambda$
$/alpha$ = $/lambda$
(b)
v
$/alpha$
$/lambda$
$/alpha$ = $/lambda$
(c)
v
$/alpha$
$/alpha$ = $/lambda$ -$/lambda$
(d)
v
$/alpha$
$/alpha$ = $/lambda$ -$/lambda$
Answer. (b)
\end{problembox}

\begin{problembox}
\textbf{Problem 69} \hfill \textbf{[Essential Physical Chemistry p.47 Q21]}
\label{q:ext-essential_all_questions-69}

When a strong acid is titrated against a strong base the end point is the point of
(a)
zero conductance
(b)
maximum conductance
(c)
minimum conductance
(d)
none of these
Answer. (c)
\end{problembox}

\begin{problembox}
\textbf{Problem 70} \hfill \textbf{[Essential Physical Chemistry p.47 Q22]}
\label{q:ext-essential_all_questions-70}

During titration of a weak acid against a weak base, there is a sharp increase in \underline{\hspace{1cm}} at the end point
(a)
conductivity
(b)
equivalent conductance
(c)
specific conductance
(d)
none of these
Answer. (a)
\end{problembox}

\begin{problembox}
\textbf{Problem 71} \hfill \textbf{[Essential Physical Chemistry p.47 Q23]}
\label{q:ext-essential_all_questions-71}

If the equivalent conductance of a certain solution of acetic acid is 39.07 $/Omega$--1 cm2 eqvt--1. If $/lambda$$/alpha$ of
CH3COOH is 390.7, the degree of dissociation of acetic acid is
(a)
0.1
(b)
0.2
(c)
0.5
(d)
0.75
Answer. (a)
\end{problembox}

\begin{problembox}
\textbf{Problem 72} \hfill \textbf{[Essential Physical Chemistry p.47 Q24]}
\label{q:ext-essential_all_questions-72}

The equivalent conductance at infinite dilution of NaCl, HCl and CH3COONa at 25$/circ$C are 126.0, 426.0
and 91.0 $/Omega$--1 cm2 respectively. The equivalent conductance of acetic acid at infinite dilution at 25$/circ$C
will be
(a)
643.0
(b)
517.0
(c)
217.0
(d)
391.0
Answer. (d)
\end{problembox}

\begin{problembox}
\textbf{Problem 73} \hfill \textbf{[Essential Physical Chemistry p.47 Q25]}
\label{q:ext-essential_all_questions-73}

The equivalent conductance of decinormal solution of acetic acid at 298 K is 80 and at infinite dilution
400 $/Omega$--1. The percentage dissociation of acetic acid is
\end{problembox}

\begin{problembox}
\textbf{Problem 74} \hfill \textbf{[Essential Physical Chemistry p.48 Q26]}
\label{q:ext-essential_all_questions-74}

If the transport number of K+ is 0.492 in KCl solution. The transport number of Cl-- ion will be
(a)
0.984
(b)
0.492
(c)
0.508
(d)
0.016
Answer. (c)
\end{problembox}

\begin{problembox}
\textbf{Problem 75} \hfill \textbf{[Essential Physical Chemistry p.48 Q27]}
\label{q:ext-essential_all_questions-75}

In a binary electrolyte AB of certain concentration, the transport number of cation is 0.4. The transport
number of the anion will be
(a)
0.2
(b)
0.4
(c)
0.6
(d)
0.8
Answer. (c)
\end{problembox}

\begin{problembox}
\textbf{Problem 76} \hfill \textbf{[Essential Physical Chemistry p.48 Q28]}
\label{q:ext-essential_all_questions-76}

The equivalent conductance at 18$/circ$C of a normal solution of KCl is 98.2 and at infinite dilution at the
same temperature is 131. The degree of dissociation of KCl will be
(a)
131/98.2
(b)
98.2/131
(c)
98.2/131 + 98.2
(d)
131/131 + 98.2
Answer. (b)
\end{problembox}

\begin{problembox}
\textbf{Problem 77} \hfill \textbf{[Essential Physical Chemistry p.48 Q29]}
\label{q:ext-essential_all_questions-77}

If the degree of dissociation of weak electrolyte at a certain temperature is 0.1, the percentage of the
molecules undissociated will be
(a)
10/%
(b)
20/%
(c)
50/%
(d)
90/%
Answer. (d)
\end{problembox}

\begin{problembox}
\textbf{Problem 78} \hfill \textbf{[Essential Physical Chemistry p.48 Q30]}
\label{q:ext-essential_all_questions-78}

For strong electrolytes, the degree of dissociation is
(a)
nearly equal to one
(b)
nearly equal to zero
(c)
nearly equal to infinity
(d)
nearly equal to 0.5
Answer. (a)
\end{problembox}

\begin{problembox}
\textbf{Problem 79} \hfill \textbf{[Essential Physical Chemistry p.48 Q31]}
\label{q:ext-essential_all_questions-79}

The equivalent conductance at infinite dilution of NH4Cl, NaOH and NaCl at 18$/circ$C are respectively,
129.8, 227.4 and 108.9 $/Omega$--1 cm2 g eqvt--1. The equivalent conductance of NH4OH at infinite dilution
at 18$/circ$C will be
(a)
128.8
(b)
108.9
(c)
227.4
(d)
238.3
Answer. (d)
\end{problembox}

\begin{problembox}
\textbf{Problem 80} \hfill \textbf{[Essential Physical Chemistry p.48 Q32]}
\label{q:ext-essential_all_questions-80}

At infinite dilution the conductance of NH4Br, KOH and KBr at 298 K are 151, 271 and 151 $/Omega$--1 cm2
g eqvt--1 respectively. If the equivalent conductance of NH4OH at 0.01 N concentration at 298 K is 27.1
$/Omega$--1 cm2 g eqvt--1, the degree of dissociation of NH4OH would be
(a)
0.1
(b)
0.01
(c)
0.001
(d)
1.0
Answer. (a)
\end{problembox}

\begin{problembox}
\textbf{Problem 81} \hfill \textbf{[Essential Physical Chemistry p.48 Q33]}
\label{q:ext-essential_all_questions-81}

$/lambda$o for Ba(OH2), BaCl2 and NH4Cl are 228, 120 and 129 $/Omega$--1 cm2 g eqvt--1. The value of $/lambda$o for NH4OH
(a)
108 $/Omega$--1 cm2 g eqvt--1
(b)
249 $/Omega$--1 cm2 g eqvt--1
(c)
348 $/Omega$--1 cm2 g eqvt--1
(d)
183 $/Omega$--1 cm2 g eqvt--1
Answer. (d)
\end{problembox}

\begin{problembox}
\textbf{Problem 82} \hfill \textbf{[Essential Physical Chemistry p.48 Q34]}
\label{q:ext-essential_all_questions-82}

The ionic mobility of an ion is given by the relation
(a)
$/lambda$$/alpha$ $/div$ 96500
(b)
$/lambda$$/alpha$ $/times$ 96500
(c)
$/lambda$$/alpha$ + 96500
(d)
$/lambda$$/alpha$ -- 96500
Answer. (a)
\end{problembox}